\chapter{The Language of Optimization Models}\label{chap:2}

\begin{quote}
\emph{``A good notation has a subtlety and suggestiveness which at times
makes it almost seem like a live teacher.''} --- Bertrand Russell
\end{quote}

This chapter introduces a common notation for discrete optimization. The
focus is on MIP- and CP-style notation that dominates much of the
modeling literature; much of the same vocabulary carries over to SAT
modeling, even though that community uses its own surface syntax. Other
communities and textbooks adopt variants---different letter conventions,
bracket styles, and shorthands for global constraints---and we do not
claim that the choices below are universal.

Fluency in this notation yields three benefits. First, it makes research
papers and textbooks readable, allowing you to extract ideas rather than
rederive them. Second, once a problem exceeds what fits in working
memory, it is faster to sketch a precise model on paper than to move
directly to code. Third, it provides a shared, unambiguous language for
personal notes, whiteboard discussions, and paper drafts.

Most of the material is standard, and the chapter also serves as a
reference. If you are unsure how much you already know, use the section
headings as a checklist.

\section{Notation in Action}\label{sec:2.1}

Recall the 0/1 knapsack: items \(I\) with weights \(w_i\) and values
\(c_i\), a knapsack of capacity \(C\), and a binary choice per item.
Compressed into a single line, the model reads

\[
\max_{x \in \{0,1\}^{|I|}} \;\Bigl\{\, \sum_{i \in I} c_i\, x_i \;\;\Big|\;\; \sum_{i \in I} w_i\, x_i \le C \,\Bigr\}.
\]

This line specifies, in order, \emph{what we choose} (a binary vector
\(x\), one component per item), \emph{what we optimize} (a weighted sum
of values), and \emph{what is allowed} (the weighted sum of weights
does not exceed \(C\)). Each symbol is functional, and a reader
familiar with the conventions can reconstruct the problem without
additional prose.

Now suppose the values are uncertain: a simulator provides \(|S|\)
scenarios, and we require a packing that performs well in the worst
case. The model changes minimally:

\[
\max_{x \in \{0,1\}^{|I|}} \; \min_{s \in S} \;\Bigl\{\, \sum_{i \in I} c^s_i\, x_i \;\;\Big|\;\; \sum_{i \in I} w_i\, x_i \le C \,\Bigr\}.
\]

The decision variables, parameters, and constraint remain unchanged;
only the objective acquires a \(\min_{s \in S}\), and the values carry a
scenario index. The formulation remains flexible: replacing
\(\min_{s \in S}\) with \(\sum_{s \in S} p_s\) yields an expected-value
objective under scenario probabilities \(p_s\); replacing it with a
quantile operator yields a value-at-risk objective. The remainder of the
model is unaffected. In contrast, a prose formulation would require
substantial rewriting for each variant.

Next, consider two knapsacks: the original with capacity \(C\), and a
second with capacity \(C'\) that only accepts light items,
\(I' = \{i \in I \mid w_i \le w'\}\). We must choose one knapsack and
pack it, while still maximizing the worst-case value. Rather than
inlining everything into a single expression, define the feasible
regions:

\[
\mathcal{F}_0 \;=\; \bigl\{\, x \in \{0,1\}^{|I|} \;\big|\; \textstyle\sum_{i \in I} w_i x_i \le C \,\bigr\},
\]

\[
\mathcal{F}_1 \;=\; \bigl\{\, x \in \{0,1\}^{|I|} \;\big|\; \textstyle\sum_{i \in I} w_i x_i \le C',\;\; x_i = 0 \;\;\forall i \in I \setminus I' \,\bigr\},
\]

and combine them using a set union:

\[
\max_{x \;\in\; \mathcal{F}_0 \,\cup\, \mathcal{F}_1} \; \min_{s \in S} \;\sum_{i \in I} c^s_i\, x_i.
\]

Each line serves a distinct purpose. The decomposition also exposes
structure: the feasible regions \(\mathcal{F}_0\) and \(\mathcal{F}_1\)
are independent of the scenarios, whereas the objective depends on them.
Thus, uncertainty affects \emph{what is optimized}, not \emph{what is
feasible}. Once defined, these regions can be reused---for example, to
state bounds, define relaxations, or compare variants---without
repeating their definitions.

The notation integrates naturally with prose: symbols appear inline
(\(I'\), \(\mathcal{F}_1\)), in displayed equations, and across
paragraphs without friction. The two registers remain consistent and
mutually reinforcing.

The remainder of this chapter develops the vocabulary underlying these
examples: sets and indices, parameters, graphs, variables and domains,
expressions, objectives, constraints, and logical connectives.

\section{Sets, Indices, and Parameters}\label{sec:2.2}

Consider a city planning problem: choose locations for ambulance
stations. Let \(C\) denote the set of neighborhoods to serve and \(F\)
the set of candidate sites. The goal is to open as few stations as
possible while ensuring that every neighborhood lies within a
five-minute drive of at least one open station. A drive-time function
\(d : C \times F \to \mathbb{R}_{\ge 0}\) returns the travel time in
minutes; for \(c \in C\) and \(f \in F\), \(d(c, f)\) is the time from
\(c\) to \(f\). For a fixed site \(f \in F\), the neighborhoods
reachable within five minutes form a subset of \(C\):

\[
N_f \;=\; \{\, c \in C \;\mid\; d(c, f) \le 5 \,\}.
\]

The braces with a vertical bar denote \emph{set-builder notation}: all
elements of \(C\) that satisfy the condition. The subscript \(f\)
indicates that this defines one set per candidate site---an indexed
family \(\{N_f\}_{f \in F}\).

\begin{platypustip}
Maintain a table of sets and parameters as a glossary. For
the ambulance problem:

\begin{symboltable}{lL}
Symbol & Meaning \\ \midrule
\(C\) & set of neighborhoods to serve \\
\(F\) & set of candidate sites \\
\(d : C \times F \to \mathbb{R}_{\ge 0}\) & drive-time function \\
\(N_f\) & neighborhoods reachable from site \(f\) within five minutes \\
\end{symboltable}


If multiple variable types appear, create an analogous table for
variables, including their domains.
\end{platypustip}

The decision is to select a subset \(F^\star \subseteq F\) of sites to
open. The coverage requirement is

\[
\bigcup_{f \in F^\star} N_f \;=\; C.
\]

Given \(F^\star\), the neighborhoods uniquely covered by an open site
\(f \in F^\star\) are

\[
N_f \;\setminus\; \bigcup_{g \in F^\star \setminus \{f\}} N_g.
\]

If this set is empty, site \(f\) is redundant. Another useful
observation: if \(N_g \subseteq N_f\) for some \(g \in F\), then \(g\)
is dominated and can be removed from \(F\) before solving, since any
coverage it provides is already achieved by \(f\).

The set operators used above, along with several standard ones:

\begin{symboltable}{lL}
\(a \in A\) & membership: \(a\) is an element of \(A\) \\
\(a \notin A\) & non-membership: \(a\) is not an element of \(A\) \\
\(A \cup B\) & union: elements in \(A\) or \(B\) (or both) \\
\(A \cap B\) & intersection: elements in both \\
\(A \setminus B\) & difference: elements in \(A\) but not in \(B\) \\
\(A \subseteq B\) & subset relation \\
\(A \times B\) & Cartesian product: ordered pairs \((a,b)\) \\
\(\lvert A \rvert\) & cardinality: number of elements \\
\(\emptyset\) & empty set \\
\(\{a \in A \mid P(a)\}\) & set-builder notation \\
\end{symboltable}

Cartesian products describe relations and graphs: for example, an arc
set \(E \subseteq V \times V\), a set of feasible (worker, job)
assignments, or the edge set of a conflict graph.

A few naming conventions recur. Sets typically use uppercase Latin
letters, with corresponding lowercase symbols for elements: \(i \in I\),
\(j \in J\), \(v \in V\). Certain letters carry conventional
meanings---\(V, E\) for vertices and edges, \(T\) for time periods,
\(K\) for constraint indices---but these are conventions, not
requirements. A symbol is defined by its declaration, not by its name.

Data can be written in two interchangeable forms. A parameter may appear
as a subscripted quantity, such as \(w_i\) or \(a_{ij}\), or as a
function, such as \(d(c, f)\). Formally, these are equivalent:
\((w_i)_{i \in I}\) and \(w : I \to \mathbb{R}\) define the same
mapping. Subscripts are compact for tabulated data; functional notation
reads more naturally when arguments span multiple index sets or when the
value is conceptually computed (e.g., a distance or cost query). The
notation \(d(c, f)\) reflects the latter interpretation.

\section{Graphs}\label{sec:2.3}

A graph is a common form of structured set: a pair \(G = (V, E)\), where
\(V\) is a finite set of vertices and \(E \subseteq V \times V\) is a
set of arcs. Vertices are typically denoted by \(v\) (or \(w\) when
needed), and \((v, w) \in E\) denotes an arc from \(v\) to \(w\), often
abbreviated as \(vw\). An \emph{undirected graph} uses unordered pairs:

\[
G = (V, E), \qquad \{v, w\} \in E.
\]

The choice between directed and undirected graphs depends on the
application---flows are directed, whereas relationships such as
friendships are not. The modeling choice can significantly affect the
resulting formulation.

For a directed graph, the \emph{out-neighborhood} and
\emph{in-neighborhood} of a vertex \(v\) are

\[
N^+(v) = \{\, w \in V \mid (v, w) \in E \,\}, \qquad N^-(v) = \{\, u \in V \mid (u, v) \in E \,\}.
\]

The notation follows the standard upstream/downstream interpretation
\(u \to v \to w\): \(u\) is a predecessor of \(v\), and \(w\) is a
successor. For an undirected graph, there is a single neighborhood,

\[
N(v) = \{\, w \in V \mid \{v, w\} \in E \,\}.
\]

A related convention reserves \(\delta^+(v)\) and \(\delta^-(v)\) for
the \emph{sets of arcs} leaving and entering \(v\),

\[
\delta^+(v) = \{\, (v, w) \in E \,\}, \qquad \delta^-(v) = \{\, (u, v) \in E \,\}.
\]

The distinction matters in flow formulations, where the natural sums are
over arcs rather than over neighboring vertices.

A \emph{subgraph} is a pair \((V', E')\) with \(V' \subseteq V\) and
\(E' \subseteq E\). It is \emph{induced} by \(V'\) if its edge set
consists exactly of those edges in \(G\) with both endpoints in \(V'\):

\[
E' = \{\, (v, w) \in E \mid v, w \in V' \,\}.
\]

The same definition applies to undirected graphs by interpreting
\((v, w)\) as \(\{v, w\}\). We write \(G[V']\) for the subgraph induced
by \(V'\). Induced subgraphs underlie structures such as cliques,
independent sets, and dense subgraphs. For example, a clique is a set
\(V' \subseteq V\) such that \(G[V']\) is complete.

As an example, consider constructing a worker's shift from a set of
tasks \(T\). Each task \(t \in T\) has a start time \(s_t\) and an end
time \(e_t \ge s_t\), and transitioning from \(t\) to \(t'\) requires
time \(\tau_{t,t'} \ge 0\). The feasible sequences of tasks can be
represented as a graph with

\[
V = T \cup \{v_{\text{start}}, v_{\text{end}}\}, \qquad
E \;=\; \bigl\{\, (t, t') \in T \times T \;\big|\; e_t + \tau_{t,t'} \le s_{t'} \,\bigr\}
\;\cup\; \bigl(\{v_{\text{start}}\} \times T\bigr)
\;\cup\; \bigl(T \times \{v_{\text{end}}\}\bigr).
\]

An arc \((t, t')\) indicates that \(t'\) can follow \(t\) in a feasible
schedule. The vertices \(v_{\text{start}}\) and \(v_{\text{end}}\)
represent the beginning and end of the shift. Since \(e_t \ge s_t\) and
\(\tau \ge 0\), all arcs move forward in time, so \(G\) is a
\emph{directed acyclic graph} (DAG).

A single worker's schedule corresponds to a path from
\(v_{\text{start}}\) to \(v_{\text{end}}\) in \(G\). If every task must
be assigned to exactly one worker and each worker performs one feasible
sequence, then minimizing the number of workers is the minimum
path-cover problem on this DAG, which reduces to bipartite matching.
With per-arc costs, the same structure yields a minimum-cost flow
formulation, typically after splitting each task vertex to enforce that
exactly one unit of flow uses it.

\begin{platypustip}
Modeling a problem as a graph provides access to a broad
range of graph algorithms, which can serve as subroutines or, in some
cases, solve the problem directly, as in this example.
\end{platypustip}

\section{Variables and Domains}\label{sec:2.4}

The decision variables \(x_i\) in the Knapsack problem lived in
\(\{0, 1\}^{|I|}\)---one bit per item, indicating whether the item is
selected. Consider the same model with a modified domain:

\[
\max \;\sum_{i \in I} c_i\, x_i
\quad\text{s.t.}\quad
\sum_{i \in I} w_i\, x_i \le C,
\qquad x_i \in [0, 1].
\]

The objective, constraint, and data are unchanged; only the domain
differs. This is the \emph{fractional knapsack} problem, where items can
be taken in any proportion between zero and one. It admits a greedy
solution in \(\mathcal{O}(|I| \log |I|)\) time by sorting items by value
density. In contrast, the original formulation with \(x_i \in \{0,1\}\)
is NP-hard. Changing the domain again to \(x_i \in \mathbb{N}_0\) yields
the \emph{unbounded knapsack} problem, where any non-negative number of
copies of each item may be selected.

Each \(x_i\) is a \emph{decision variable}---a quantity determined by
the solver---and its \emph{domain} specifies the admissible values.
Common domains include:

\begin{symboltable}{lL}
Domain & Typical use \\ \midrule
\(\mathbb{R}\) & continuous quantities (flows, prices, fractions) \\
\(\mathbb{Z}\) & integers (indices, offsets) \\
\(\mathbb{N}_0\) & non-negative integers (counts, quantities) \\
\(\{0,1\}\) & binary decisions \\
\(\{1,2,\dots,k\}\) & finite enumerations (machines, modes) \\
\(\{\text{red}, \text{green}\}\) & symbolic enumerations \\
\end{symboltable}


Bounds are a special case of domains: \(l_i \le x_i \le u_i\) is
equivalent to \(x_i \in [l_i, u_i]\), and \(x_i \ge 0\) to
\(x_i \in [0, \infty)\). When multiple variables share a domain, we
write \(x \in \mathbb{R}^n\), \(x \in \mathbb{Z}^n\), or
\(x \in \{0,1\}^n\).

Variable names follow informal conventions---\(x\) for primary
decisions, \(y\) for auxiliary variables, \(s\) for slack variables, and
\(z\) for objective values---but these conventions are not binding.

\section{Expressions}\label{sec:2.5}

Expressions are quantities formed from variables and parameters; once
the variables are fixed, an expression evaluates to a scalar. It is
useful to distinguish between linear and nonlinear expressions, because
the distinction has algorithmic consequences. Models composed entirely
of linear expressions---\emph{linear programs} and \emph{mixed-integer
linear programs}---are significantly better understood than their
nonlinear counterparts. In practice, nonlinear formulations are often
transformed into linear ones, either exactly or through approximation.

\subsection{Linear expressions}\label{sec:2.5.1}

A linear expression is a weighted sum of variables:

\[
\sum_{i \in I} c_i x_i.
\]

This form appears in many contexts: total value in knapsack, total
assignment cost, or total flow into a vertex. Nested sums extend the
structure over Cartesian products:

\[
\sum_{i \in I} \sum_{j \in J} a_{ij} x_{ij}
\;=\;
\sum_{(i,j) \in I \times J} a_{ij} x_{ij}.
\]

The right-hand form is convenient when the index set is a subset of
\(I \times J\), such as the edge set of a graph:

\[
\sum_{(i,j) \in E} x_{ij}.
\]

In vector notation, \(\sum_i c_i x_i\) is written as \(c^\top x\), and a
system of linear constraints as \(A x \le b\). The indexed and vector
forms represent the same object, differing only in presentation.

\subsection{Nonlinear expressions}\label{sec:2.5.2}

Any expression that is not a linear combination of variables is
\emph{nonlinear}. Three common patterns recur.

The first consists of variables passed through nonlinear functions, such
as products or transcendental functions:

\[
x_i^2, \quad x_i x_j, \quad \exp(x_i), \quad \log(x_i).
\]

The second pattern aggregates over an index set, analogous to \(\sum\)
in \Cref{sec:2.5.1}, but produces a nonlinear result:

\[
\max_{i \in I} g_i(x), \quad \min_{i \in I} g_i(x), \quad
\operatorname{median}_{i \in I} g_i(x), \quad
\operatorname{mode}_{i \in I} g_i(x), \quad
\operatorname{quantile}_{i \in I}^{\,q} g_i(x).
\]

The third pattern connects logical predicates to numeric expressions.
The \emph{Iverson bracket} maps a predicate \(P\) to a binary value:

\[
[P] \;=\; \begin{cases} 1 & \text{if } P \text{ is true}, \\ 0 & \text{otherwise.} \end{cases}
\]

For example, the number of variables taking a specific value is

\[
\sum_{i \in I} [\, x_i = v \,],
\]

and the condition ``at most \(k\) of the predicates \(P_1, \dots, P_n\)
hold'' can be written as \(\sum_{i=1}^{n} [P_i] \le k\).

\subsection{Declaring functions}\label{sec:2.5.3}

A function is declared by specifying its \emph{signature}:

\[
f : A \to B,
\]

which defines the function name \(f\), its domain \(A\), and its
codomain \(B\). Optionally, a definition follows:

\[
f(a) := \dots.
\]

The domain and codomain may be Cartesian products, as in
\(d : C \times F \to \mathbb{R}_{\ge 0}\). This notation covers
parameters (data expressed as functions rather than indexed symbols),
derived expressions (e.g., \(\mathrm{slack}(x) := C - \sum_i w_i x_i\)),
and predicates (with codomain \(\{0,1\}\) or
\(\{\text{true}, \text{false}\}\)).

\begin{platypusinfo}
An optimization problem can itself be viewed as a
function: parameters map to an optimal value. For the 0/1 knapsack,
\begin{align*}
\mathrm{knapsack} :\; & \mathbb{R}_{\ge 0}^{n} \times \mathbb{R}_{\ge 0}^{n} \times \mathbb{R}_{\ge 0} \to \mathbb{R}_{\ge 0}, \\
\mathrm{knapsack}(c, w, C) \;=\; & \max_{x \in \{0,1\}^{n}} \Bigl\{\, \textstyle\sum_{i} c_i x_i \;\Big|\; \sum_i w_i x_i \le C \,\Bigr\}.
\end{align*}
\end{platypusinfo}

\section{Objectives}\label{sec:2.6}

An optimization problem selects \(x\) from a feasible region
\(\mathcal{F}\) to minimize a scalar objective \(f : D \to \mathbb{R}\):

\[
\min_{x \in \mathcal{F}} f(x).
\]

Maximization replaces \(\min\) with \(\max\), or equivalently
minimizes \(-f\). The minimum attained is the \emph{optimal value},

\[
f^\star = \min_{x \in \mathcal{F}} f(x),
\]

and the set of \emph{optimal solutions} is

\[
\arg\min_{x \in \mathcal{F}} f(x) \;=\; \{\, x \in \mathcal{F} \mid f(x) = f^\star \,\}.
\]

We write \(x^\star \in \arg\min\) rather than \(x^\star = \arg\min\),
because the optimum need not be unique; many problems admit multiple
optimal solutions, sometimes forming a continuum.

When multiple criteria are considered, we obtain a vector of objective
values \((f_1(x), \dots, f_m(x))\). Such vectors are not minimized by
the usual scalar order; one must choose a notion of preference, such as
Pareto dominance, a weighted sum, a minimax objective, or a
lexicographic order. The choice is a modeling decision and is not
pursued here.

\section{Constraints}\label{sec:2.7}

A constraint is a \emph{predicate}---a statement that evaluates to true
or false for a given assignment of values. While expressions answer
``what is the value?'', constraints answer ``is the condition
satisfied?''.

The most common predicates in optimization are linear (in)equalities:

\[
\sum_{j \in J} a_{ij} x_j \;\le\; b_i \qquad \forall i \in I.
\]

The quantifier \(\forall i \in I\) is essential: this defines a family
of \(|I|\) constraints, not a single one. For example, if
\(I = \{1,2,3\}\), the family expands to

\[
\sum_j a_{1j} x_j \le b_1, \qquad \sum_j a_{2j} x_j \le b_2, \qquad \sum_j a_{3j} x_j \le b_3.
\]

Equalities and reverse inequalities follow the same pattern.

Linearity is convenient but not required. In general, a constraint has
the form

\[
P_k(x) \qquad \forall k \in K,
\]

where each \(P_k\) is a predicate on the variable vector---for example,
a linear inequality, an \emph{AllDifferent} constraint, a
\emph{NoOverlap} constraint, or another relation.

\section{Logic and Quantifiers}\label{sec:2.8}

Many decisions are yes/no decisions: \emph{open this station},
\emph{take this course}, \emph{make this assignment}. A 0/1 variable is
enough to encode them. Models built from such Boolean variables use the
standard logical connectives

\[
\land, \quad \lor, \quad \neg, \quad \Rightarrow, \quad \Leftrightarrow,
\]

possibly quantified over index sets with \(\forall\) and \(\exists\).
Two common large-operator forms are

\[
\bigwedge_{k \in K} \phi_k, \qquad \bigvee_{k \in K} \phi_k.
\]

They denote the conjunction and disjunction of an indexed family of
Boolean expressions \(\phi_k\). The first is equivalent to the
constraint family \(\phi_k ;\forall k \in K\); the second means that at
least one \(\phi_k\) must hold, and therefore cannot be replaced by a
family of separate constraints.

\begin{center}\rule{0.5\linewidth}{0.5pt}\end{center}

\paragraph{Example.} After an office prank, three suspects remain: Anna,
Ben, and Cara. Let \(x_a, x_b, x_c \in \{0,1\}\) encode whether each
person is guilty. Witnesses provide four facts. Anna would only act with
Ben, so \(x_a \Rightarrow x_b\). Ben and Cara have never gotten along,
so \(\neg(x_b \land x_c)\). Cara would only join if Anna did, so
\(x_c \Rightarrow x_a\). Finally, someone did it, so
\(x_a \lor x_b \lor x_c\). Together, the facts give

\[
(x_a \Rightarrow x_b) \;\land\; \neg(x_b \land x_c) \;\land\; (x_c \Rightarrow x_a) \;\land\; (x_a \lor x_b \lor x_c).
\]

Chaining the implications gives \(x_c \Rightarrow x_a \Rightarrow x_b\).
But \(x_b\) implies \(\neg x_c\), because Ben and Cara cannot both be
guilty. Hence \(x_c \Rightarrow \neg x_c\), so \(x_c = 0\). With Cara
cleared, the condition that someone is guilty reduces to
\(x_a \lor x_b\). Together with \(x_a \Rightarrow x_b\), this forces
\(x_b = 1\), regardless of Anna's role. The clues implicate Ben, clear
Cara, and leave Anna's status undetermined.

\begin{center}\rule{0.5\linewidth}{0.5pt}\end{center}

The connectives are not limited to Boolean variables; they also compose
predicates over richer variables. Suppose two jobs share a machine. With
start times \(s_1, s_2 \in \mathbb{N}_0\) and durations \(d_1, d_2\),
the jobs must run one after the other:

\[
s_1 + d_1 \le s_2 \;\;\lor\;\; s_2 + d_2 \le s_1.
\]

This is a disjunction of two linear inequalities over integer variables.
The same logical connectives apply to inequalities, equalities,
\(\text{AllDifferent}(y_1,\ldots,y_k)\), \(\text{NoOverlap}\)
constraints, or any other relation. A predicate, like a function, can be
named for reuse:

\[
P(x) \;:=\; \sum_{i \in I} w_i x_i \le C.
\]

Once named, it fits into the same logical syntax:
\(P_1(x) \land P_2(x)\), \(P(x) \Rightarrow Q(x)\), or
\(\bigvee_{k \in K} P_k(x)\).

The two layers---Boolean variables and predicates over richer
variables---meet in \emph{reification}. Reification ties a Boolean
indicator \(z \in \{0,1\}\) to the truth value of a predicate:

\[
z \;\Leftrightarrow\; (x \ge 5).
\]

Once reified, \(z\) can appear in logical connectives like any other
Boolean variable. The special case \(z \Rightarrow a^\top x \le b\) is
known as an \emph{indicator constraint} in MIP terminology. How a solver
encodes such patterns is separate from how the model states them.

\section{A Standard Model Template}\label{sec:2.9}

Many optimization models share a common skeleton. In its simplest form,
it is

\[
\begin{aligned}
\min_{x} \quad      & f(x) \\
\text{s.t.} \quad   & P_k(x) && \forall k \in K, \\
                    & x \in D.
\end{aligned}
\]

The objective \(f\), predicates \(P_k\), and domain \(D\) are specified
by the application. The corresponding feasible region is

\[
\mathcal{F} = \{\, x \in D \mid P_k(x) \;\forall k \in K \,\},
\]

with optimal value and optimal solution written as

\[
f^\star = \min_{x \in \mathcal{F}} f(x),\qquad
x^\star \in \arg\min_{x \in \mathcal{F}} f(x).
\]

Many formulations you encounter will fit this template, often with only
minor notational variations.
